\section{Conclusions}
This paper presented ADAPT, an adaptive dual-projection framework for perceptive humanoid traversal in complex 3D environments. By decomposing local geometry into a horizontal elevation map and a vertical distance map, ADAPT explicitly separates ground-support and body-level clearance constraints into a compact representation for high-frequency control. Furthermore, embedding the sensing radius into the action space enables the robot to dynamically adjust its perceptual horizon based on motion context. Experiments demonstrate that ADAPT significantly improves traversal robustness and computational efficiency over fixed-range and voxel-based alternatives, achieving successful zero-shot sim-to-real transfer on a Unitree G1 robot.

Despite these capabilities, the current framework focuses solely on geometric traversability and adapts only the sensing radius. Future work will investigate richer active perception policies that jointly optimize sensing scale, resolution, and update frequency, while incorporating semantic reasoning to enable long-horizon navigation in complex, human-centered environments.